% Generated draft: requires mathematical and source review.
\documentclass[11pt]{article}
\usepackage[margin=25mm]{geometry}
\usepackage{fontspec}
\setmainfont{Latin Modern Roman}
\usepackage{amsmath,amssymb,mathtools,unicode-math}
\setmathfont{Latin Modern Math}
\usepackage{enumitem,needspace,xurl,hyperref}
\setlength{\parindent}{0pt}
\setlength{\parskip}{5pt}
\setlength{\emergencystretch}{3em}
\newcommand{\N}{\mathbb{N}}
\newcommand{\Z}{\mathbb{Z}}
\newcommand{\Q}{\mathbb{Q}}
\newcommand{\R}{\mathbb{R}}
\newcommand{\C}{\mathbb{C}}
\newcommand{\F}{\mathbb{F}}
\newcommand{\A}{\mathbb{A}}
\newcommand{\PP}{\mathbb{P}}
\newcommand{\E}{\mathbb{E}}
\newcommand{\eps}{\varepsilon}
\newcommand{\dd}{\,\mathrm{d}}
\begin{document}
\section*{UnsolvedMath ID: 7200012}
\textbf{Source problem code:} catalogue entry
\subsection{Definitions and mathematical statement}
The record names the Maulik--Nekrasov--Okounkov--Pandharipande correspondence between Gromov--Witten and Donaldson--Thomas theory, without giving a precise version. A standard unmarked Calabi--Yau threefold formulation illustrates the required data, rather than defining every version of the named conjecture.

Let $X$ be a smooth projective complex threefold with trivial canonical bundle. For a nonzero effective curve class $\beta$, let $N_{g,\beta}$ be the zero-point genus-$g$ Gromov--Witten invariant, the degree of the virtual class of stable maps in class $\beta$. Let $I_{n,\beta}$ be the Donaldson--Thomas invariant of ideal sheaves of subschemes $Z\subset X$ with one-dimensional class $[Z]=\beta$ and Euler characteristic $\chi(\mathcal O_Z)=n$. These virtual classes have dimension zero in this setting. With $Q^\beta$ formal curve-class variables, form
\[
Z'_{\mathrm{GW}}(u,Q)=\exp\left(\sum_{\beta\ne0}\sum_{g\geq0}
N_{g,\beta}u^{2g-2}Q^\beta\right),
\]
\[
Z'_{\mathrm{DT}}(q,Q)=
\frac{\sum_{\beta,n}I_{n,\beta}q^nQ^\beta}
{\sum_n I_{n,0}q^n}.
\]
The first expression omits degree-zero contributions; the second divides them out. The conventional correspondence requires appropriate rationality in $q$ of fixed-curve-class reduced DT series and identifies their expansion under
\[
q=-e^{iu},\qquad Z'_{\mathrm{GW}}(u,Q)=Z'_{\mathrm{DT}}(-e^{iu},Q).
\]
The change of variable uses the rational functions and their prescribed expansion, not an unjustified substitution into an arbitrary $q$-adic series. Curve variables use a suitable formal/Novikov completion; no analytic convergence is asserted.

This is one usual precise candidate suggested by the summary's Calabi--Yau-threefold context. The source does not specify projectivity, smoothness, insertions/descendants, equivariant or noncompact variants, connected versus disconnected conventions, rationality, degree-zero normalization, or the variable change. General threefold and inserted correspondences require different data; they are not automatically covered by the displayed unmarked Calabi--Yau version. A source-faithful full formalization must first identify the intended version.
\subsection{Short English statement}
The MNOP conjecture relates suitably normalized curve-counting partition functions in Gromov--Witten and ideal-sheaf Donaldson--Thomas theory after a change of variable. The record gives only its name, so the displayed Calabi--Yau version is an explicit conventional interpretation requiring primary-source review.
\subsection{Sources}
\begin{itemize}
\item Local UnsolvedMath registry, permanent ID 7200012, exact source code catalogue entry, named for Maulik--Nekrasov--Okounkov--Pandharipande; no proposer metadata or year supplied. Registry status partially solved, classification PARTIAL-PROGRESS.
\item Wikipedia contributors, \emph{List of unsolved problems in mathematics}, Geometry bullet 12, pinned revision 1366636767, \url{https://en.wikipedia.org/w/index.php?title=List_of_unsolved_problems_in_mathematics&oldid=1366636767#Geometry}. Supplied MediaWiki source extraction accessed July 29, 2026. No primary MNOP paper or exact version is cited in the record.
\item The supplied research summary reports the correspondence established in toric and many geometric families while describing arbitrary Calabi--Yau threefold generality as open as of 2026. This is a supplied unverified assessment, not a newly performed literature check or a theorem with specified hypotheses.
\item Provenance says the extracted bullet remains unpublished pending primary-source status review, although registry metadata has published set true; this administrative discrepancy is preserved. The pinned list is not a primary resolution source.
\item Wikipedia text is CC BY-SA 4.0; attribution and share-alike handling remain required for release. Status is flagged NEEDS\_REVIEW.
\end{itemize}
\end{document}
